\documentclass[10pt,a4paper]{amsart}
\usepackage[margin=19mm]{geometry}
\usepackage{amsmath,amssymb,mathtools}
\usepackage[scr=rsfs,cal=euler]{mathalfa}
\usepackage{xfrac}
\usepackage[unicode,hidelinks,bookmarks=false]{hyperref}
\newcommand{\ie}{i.e.\ }
\newcommand{\cf}{cf.\ }
\newcommand{\resp}{respectively\ }
\newcommand{\Subsection}{\subsection}
\providecommand{\nobreakdash}{\nobreak\hbox{-}}
\setlength{\emergencystretch}{2em}
\multlinegap0pt
\usepackage{graphics}
\usepackage{epsfig}
\usepackage{amscd}
\numberwithin{equation}{section}
\usepackage{mathrsfs}
\usepackage{color}
\usepackage{bm}
\usepackage{float}
\usepackage{tikz}
\usepackage{multirow}
\usepackage{verbatim}
\usepackage{placeins}
\usetikzlibrary{matrix,arrows,arrows.meta,calc}
\usetikzlibrary{decorations.pathreplacing,decorations.markings}
\usetikzlibrary{shapes,positioning}
\tikzstyle arrowstyle=[scale=1]
\tikzstyle directed=[postaction={decorate,decoration={markings,
    mark=at position .5 with {\arrow[arrowstyle]{stealth}}}}]
\tikzstyle reverse directed=[postaction={decorate,decoration={markings,
    mark=at position .65 with {\arrowreversed[arrowstyle]{stealth};}}}]
\newtheorem{thm}{Theorem}[section]
\newtheorem{cor}[thm]{Corollary}
\newtheorem{prop}[thm]{Proposition}
\newtheorem{conj}[thm]{Conjecture}
\newtheorem{prob}[thm]{Problem}
\newtheorem{lem}[thm]{Lemma}
\newtheorem{rem}[thm]{Remark}
\newtheorem{defn}[thm]{Definition}
\newtheorem{ex}[thm]{Example}
\numberwithin{equation}{section}
\newcommand{\C}{\mathbb{C}}
\newcommand{\Q}{\mathbb{Q}}
\newcommand{\R}{\mathbb{R}}
\newcommand{\Z}{\mathbb{Z}}
\newcommand{\Frac}{\mathrm{Frac}}
\newcommand{\A}{\mathcal{A}}
\newcommand{\B}{\mathcal{B}}
\newcommand{\F}{\mathcal{F}}
\newcommand{\FF}{\mathscr{F}}
\newcommand{\Flag}{\mathrm{Flag}}
\newcommand{\T}{\mathcal{T}}
\newcommand{\M}{\mathcal{M}}
\newcommand{\Ch}{\mathrm{Ch}}
\newcommand{\V}{\mathcal{V}}
\newcommand{\VG}{\mathrm{VG}}
\newcommand{\sgn}{\mathrm{sgn}}
\newcommand{\supp}{\mathrm{supp}}
\DeclareMathOperator{\rank}{rank}
\DeclareMathOperator{\diam}{diam}
\DeclareMathOperator{\codim}{codim}
\begin{document}
\title[Cremona invariance of filtered Varchenko--Gelfand algebras]{Cremona invariance of filtered Varchenko--Gelfand algebras}
\author{Ye Liu}
\date{}
\maketitle
\noindent\emph{Source and licence.} Cremona invariance of filtered Varchenko--Gelfand algebras. Version arXiv:2607.15787v1 (CC BY 4.0). This material is distributed under the Creative Commons Attribution 4.0 International licence, http://creativecommons.org/licenses/by/4.0/, as recorded in the arXiv OAI licence metadata for this exact version and shown on the abstract page https://arxiv.org/abs/2607.15787v1. Full rights notice: licenses/arxiv-2607.15787-CC-BY-4.0.txt.\par\smallskip
\noindent\textit{Editorial scope.} Counted unit: the original \texttt{thm} environment \texttt{main} at source lines 173--187 together with its complete proof at lines 205--268; the delivered document keeps source lines 160--269 so that every referenced label and every citation resolves inside it. Native editable source is the paper's own concatenated TeX; the AMS wrapper is editorial formatting only. No publisher class, upstream script or unrelated artwork is redistributed.
\section{Results}\label{cremona}

Suppose that $\A$ is an arrangement in $\R^r$ with coordinates $x_1,\ldots,x_r$ containing all the coordinate hyperplanes $x_i=0$, and that every other defining form is supported on two coordinates:
\[
\alpha_e=a_ex_{i(e)}+b_ex_{j(e)},
\qquad i(e)\neq j(e),\qquad a_eb_e\neq0.
\]
We define the \emph{Cremona transform} $\A^\kappa$ as an arrangement in $\R^r$ with coordinates $X_1,\ldots,X_r$, by retaining the coordinate hyperplanes $X_i=0$ and replacing $\alpha_e$ with
\[
\beta_e=b_eX_{i(e)}+a_eX_{j(e)}.
\]
This operation is induced on the complements by coordinatewise inversion.


\begin{thm}\label{main}
For every commutative ring $R$, coordinatewise inversion
\[
\kappa(x_1,\ldots,x_r)=(X_1,\ldots,X_r)=\left(x_1^{-1},\ldots,x_r^{-1}\right)
\]
induces an isomorphism of filtered $R$-algebras
\[
\kappa^*:\VG(\A^\kappa)_R\xrightarrow{\ \cong\ }\VG(\A)_R.
\]
Equivalently,
\[
\kappa^*F^p\VG(\A^\kappa)_R=F^p\VG(\A)_R
\qquad\text{for every }p\geq0.
\]
\end{thm}



We first record the elementary rank-two identity underlying the construction.  For a nonzero real number $u$, write $s_u=\sgn(u)\in\{\pm1\}$ and $h_u=(s_u+1)/2\in\{0,1\}$.

\begin{lem}\label{sign-identity}
If $u$, $v$, and $u+v$ are nonzero real numbers, then
\[
s_us_vs_{u+v}=s_u+s_v-s_{u+v}.
\]
\end{lem}

\begin{proof}
The sign patterns $(+1,+1,-1)$ and $(-1,-1,+1)$ cannot occur for $(s_u,s_v,s_{u+v})$.  The identity is immediate on each of the six remaining sign patterns.
\end{proof}

We now prove the main theorem.


\begin{proof}[Proof of Theorem \ref{main}]
Because $\A$ and $\A^\kappa$ contain all coordinate hyperplanes, both complements lie in the real torus $(\R^\times)^r$.  For a noncoordinate defining form of $\A$
\[
\alpha_e=a_ex_i+b_ex_j
\]
and its transformed form in $\A^\kappa$
\[
\beta_e=b_eX_i+a_eX_j,
\]
we have
\begin{equation}\label{pullback}
\kappa^*\beta_e
=\frac{b_e}{x_i}+\frac{a_e}{x_j}
=\frac{a_ex_i+b_ex_j}{x_ix_j}
=\frac{\alpha_e}{x_ix_j}.
\end{equation}
Thus $\kappa$ restricts to a homeomorphism of complements
\[
\kappa:M(\A)\xrightarrow{\ \cong\ }M(\A^\kappa)
\]
and induces a bijection of chamber sets $\kappa: \Ch(\A)\to\Ch(\A^\kappa)$.  Pullback along this chamber bijection is therefore an isomorphism of the underlying function algebras.

It remains to check the filtration.  Put
\[
\varepsilon_e=\sgn(a_e),\qquad \delta_e=\sgn(b_e).
\]
Equation \eqref{pullback} gives
\[
\kappa^*s_{\beta_e}=s_{x_i}s_{x_j}s_{\alpha_e}.
\]
Apply Lemma \ref{sign-identity} to $u=a_ex_i$ and $v=b_ex_j$.  Since
\[
s_u=\varepsilon_es_{x_i},\qquad
s_v=\delta_es_{x_j},\qquad
s_{u+v}=s_{\alpha_e},
\]
we obtain
\begin{equation}\label{sign-pullback}
\kappa^*s_{\beta_e}
=\delta_es_{x_i}+\varepsilon_es_{x_j}
-\varepsilon_e\delta_es_{\alpha_e}.
\end{equation}
Substituting $s=2h-1$ and simplifying as an identity of $\Z$-valued functions yields
\begin{equation}\label{heaviside-pullback}
\kappa^*h_{\beta_e}
=\delta_eh_{x_i}+\varepsilon_eh_{x_j}
-\varepsilon_e\delta_eh_{\alpha_e}
+\frac{(1-\varepsilon_e)(1-\delta_e)}{2}.
\end{equation}
The final constant is either $0$ or $2$, so the right-hand side is an integral affine-linear combination of Heaviside functions.  For the coordinate generators one simply has
\[
\kappa^*h_{X_i}=h_{x_i}.
\]
Consequently,
\[
\kappa^*F^1\VG(\A^\kappa)_R\subseteq F^1\VG(\A)_R
\]
for every commutative ring $R$.  Because the filtration is multiplicatively generated in degree one, it follows that
\[
\kappa^*F^p\VG(\A^\kappa)_R\subseteq F^p\VG(\A)_R
\qquad\text{for all }p\geq0.
\]
Finally, $(\A^\kappa)^\kappa=\A$ and $\kappa^2=\mathrm{id}$ on the torus.  Applying the same argument in the opposite direction gives the reverse inclusion.  Hence every filtered piece is preserved.
\end{proof}


\end{document}
